\chapter{Kolmogorov-Komplexität}\label{ch:kolmogorov}
Wir führen hier den wichtigen Begriff der Kolmogorov-Komplexität ein. Die Kolmogorov-Komplexität gibt uns eine Möglichkeit den Informationsinhalt von Texten zu bestimmen und erlaubt uns über kürzeste Darstellungen zu sprechen.
Die Kolmogorov-Komplexität gibt uns auch eine sinnvolle Definition der Bedeutung von Zufälligkeit von Texten.

\section{Intuitiion}
In diesem Kapitel fassen wir Wörter (Texte) als von Information auf.
Wir versuchen eine robuste Methode zur Messung des Informationsgehalts eines Wortes zu entwickeln.

\begin{myexample}\label{myexample:passwort}
Alice hat ein langes binäres Passwort für den Zugriff auf eine Streaming-Platform gewählt.
Bob möchte gerne die neueste Staffel seiner Lieblingsserie schauen und ruft dazu Alice an um sie um das Passwort zu bitten\footnote{Bob scheint nicht mit dem \textit{BitTorrent} Kommunikationsprotokoll (peer-to-peer file sharing) vertraut zu sein.
	Vielleicht benötigt er aber auch nur einen Vorwand, um Alice anzurufen.}.

\begin{itemize}
	\item Angenommen Alices Passwort lautet $p_0 := 01011010110001101101111010$. Wie könnte Sie Bob das Passwort $p_0$ per Telefon diktieren? Ihr wird (mehr oder weniger) nichts Schlaueres übrig bleiben, als das Passwort Bob Bit für Bit (Ziffer um Ziffer) zu diktieren.
	\item Anders hingegen wäre die Situation bei dem Passwort
	      \begin{align*}
		      p_1 := 11100111001110011100111001110011100 = (11100)^7,
	      \end{align*}
	      Das Passwort $p_1$ weist eine starke Regelmässigkeit auf. Natürlich könnte Alice auch dieses Passwort Bit für Bit an Bob diktieren. Offensichtlich ist es jedoch einfacher, wenn Sie ihm einfach sagen würde:
	      \begin{center}
		      \enquote{Du erhältst das Passwort, wenn du siebenmal das Wort $11100$ hintereinander schreibst.}
	      \end{center}
	      Dadurch hat Alice die Darstellung von $p_1$ \textit{komprimiert}.
\end{itemize}
\end{myexample}

\cref{myexample:passwort} motiviert die folgende intuitive Vorstellung: Wir wollen einem Wort einen kleinen Informationsgehalt beimessen, falls es eine kurze Darstellung besitzt (komprimierbar ist) und
einen grossen Informationsgehalt, falls das Wort so unregelmässig ist, dass keine kurze Darstellung zulässt.

\begin{mydefinition}[Komprimierung]\label{mydefinition:Komprimierung}
Die Erstellung einer kürzeren Darstellung eines Wortes $w$ nennen wir eine
Komprimierung von $w$.
\end{mydefinition}

\begin{myremark}
	Wir werden in diesem Kapitel nur binäre Wörter betrachten.
\end{myremark}

Ein denkbares Vorgehen zur Messung des Informationsgehalts eines binären Worts $w$ wäre, sich auf eine konkrete Komprimierungsmethode $komp$ zu einigen.
Die Länge $\abs{komp(w)}$ des komprimierten Wortes $komp(w)$ könnte dann als Mass für den Informationsgehalt von $w$ angesehen werden.

Natürlich muss auch $komp(w)$ ein Wort über dem binären Alphabet sein.
Würden wir $komp(w)$ in einem anderen Alphabet darstellen, wäre der Vergleich der Darstellungslängen nicht mehr sinnvoll, denn schliesslich ist es immer möglich, ein Wort in einem mächtigen Alphabet kurz darzustellen.

\section{Problematik einer fest gewählten Komprimierung}
Es existieren unendlich viele Komprimierungsmethoden, welche für eine feste Wahl einer Komprimierung zur Verfügung stehen.
Doch welche Komprimierungsmethode stellt nun die korrekte Wahl dar? Das Problem wird sein, dass, egal welche Komprimierung verwendet wird, immer eine andere Komprimierung existiert, welche für unendlich viele Wörter kürzere Darstellungen erzeugt.

\begin{myexample}[Komprimierung durch Ausnutzung von Wiederholungen]\label{myexample:Wiederholungen}
Lassen Sie uns ein binäres Wort mit zahlreichen Wiederholungen von Teilwörtern untersuchen. Konkret wollen wir das Wort    \begin{align*}
	w := 00(101)^{\textcolor{Red}{25}}(01010)^{\textcolor{Blue}{9}}(1011)^{\textcolor{Green}{16}} = 00(101)\textcolor{Red}{11001}(01010)\textcolor{Blue}{1001}(1011)\textcolor{Green}{10000}
\end{align*}
untersuchen, wobei wir die Exponenten in Basis 10 durch ihre binäre Darstellung ersetzt haben. Die resultierende Darstellung ist jedoch noch immer kein binäres Wort. Tatsächlich ist es ein Wort über dem Alphabet $\lrc{0, 1, (, )}$, welches vier Symbole erhält. Mit der Vereinbarung
\begin{align*}
	0 \to 00, \quad 1 \to 11, \quad ( \to 10, \quad ) \to 01
\end{align*}
erhalten wir die doppelt so lange (aber dafür binäre) Darstellung:
\begin{align*}
	w = 00001011001101111100001110001100110001110000111011001111011100000000.
\end{align*}
\end{myexample}

\begin{myexample}[Komprimierung durch Primfaktorisierung]\label{myexample:Primkomprimierung}
Es sei $x$ ein binäres Wort. Falls $\text{Nummer}(x) \geq 2$, dann gilt
\begin{align*}
	\text{Nummer}(x) = p_1^{e_1}p_2^{e_2}\ldots p_t^{e_t},
\end{align*}
wobei die $p_i$ unterschiedliche Primzahlen sind und $e_i\geq 0$ für $i\in\setcm{n\in\N}{1\leq n\leq t}$.
Eine denkbare Darstellung von $p_1^{e_1}p_2^{e_2}\ldots p_t^{e_t}$ ist
\begin{align*}
	\text{Bin}(p_1)(\text{Bin}(e_1))\text{Bin}(p_2)(\text{Bin}(e_2))\ldots \text{Bin}(p_t)(\text{Bin}(e_t)).
\end{align*}
Mit der Vereinbarung aus \cref{myexample:Wiederholungen} erhalten wir wieder eine binäre Darstellung von $x$.
\end{myexample}
Leider sind die beiden Komprimierungsmethoden aus \cref{myexample:Wiederholungen,myexample:Primkomprimierung} nicht miteinander vergleichbar. Die erste Methode (Ausnutzung von Wiederholungen) erzeugt für gewisse Wörter eine kürzere Darstellung als die zweite Methode (Primfaktorisierung) und umgekehrt.

Die Definition eines Komplexitätsmasses sollte aber robust sein in dem Sinne, dass sie nicht auf einer arbiträren Wahl einer Komprimierungsmethode basiert. Ein Komplexitätsmass, welche keine arbiträre Wahl verwendet ist die Kolmogorov-Komplexität. Um deren Definition zu verstehen, müssen wir zunächst den Begriff der Generierung eines Worts einführen.

\begin{mydefinition}[Generierung eines Worts]\label{mydefinition:erzeugung}
Es sei $w$ ein beliebiges Wort. Wir sagen, dass ein Programm (Algorithmus) $A$ das Wort $w$ generiert, falls der Aufruf $A()$ das Wort $w$ ausgibt.
\end{mydefinition}

\begin{myexample}
Das Programm
\begin{lstlisting}
        def A():
            print("00110111")
    \end{lstlisting}
generiert das Wort $00110111$.
\end{myexample}

\begin{mydefinition}[Kolmogorov-Komplexität]\label{mydefinition:kolmogorov}
Es sei $w$ ein binäres Wort. Die Kolmogorov-Komplexität $K(w)$ von $w$ ist definiert als das Minimum der \textbf{binären Längen} aller Python-Programme, die $w$ generieren.
\end{mydefinition}
Für ein binäres Wort $w$ betrachten wir alle (unendlich vielen) Maschinencodes von Python-Programmen, die $w$ generieren. Die Länge eine kürzesten\footnote{Im Allgemeinen kann es mehrere verschiedene kürzeste Maschinencodes geben, die $w$ generieren.} solcher Maschinencodes ist dann $K(w)$.

\begin{myremark}
	Ist $K(w)$ nun ein geeignetes Mass für den Informationsgehalt des binären Wortes $w$? Ja, da jede Komprimierungsmethode als Programm formuliert werden kann, bezieht die Kolmogorov-Komplexität jede Komprimierungsmethode ein.
	Es sei $x$ ein binäres Wort und $f$ eine Komprimierungsmethode, die zu $x$ eine komprimierte Darstellung $x' := f(x)$ generiert. Es sei $f^{-1}$ die Umkehrung der Komprimierung $f$, also $f^{-1}(f(x')) = x$. Dann können wir ein Programm schreiben, welches (das kurze Wort) $x'$ als Parameter beinhaltet und $x$ wieder aus $x'$ mithilfe von $f^{-1}$ erzeugt und schliesslich ausgibt. Dazu muss $f^{-1}$ natürlich auch im Programm kodiert sein. Doch die binäre Länge der Beschreibung dieser Umkehrfunktion ist unabhängig von der Längen von $x$ und $x'$ und kann somit als konstant angesehen werden.
\end{myremark}

\begin{mytheorem}[$K(w)$ ist auf keiinen Fall wesentlich länger als $\abs{w}$]
	Es existiert eine Konstante $c$, so dass für jedes binäre Wort $w$ gilt
	\begin{align*}
		K(w) \leq \abs{w} + c.
	\end{align*}
	\begin{myproof}
		Es genügt, ein Programm anzugeben, welches $w$ generiert und eine binäre Länge aufweist, welche kleiner oder gleich $\abs{w} + c$ ist (für eine geeignete Konstante $c$). Das Programm
		\begin{lstlisting}
            def A():
                print(x)
        \end{lstlisting}
		generiert $x$. Doch wie gross ist die binäre Länge von $A$?
	\end{myproof}
\end{mytheorem}
